\chapter{Transformación de Coordenadas}

\section{Motivación}

Al estudiar un sistema mecánico, la elección del sistema de coordenadas no
es única: un mismo punto del espacio se puede describir en coordenadas
cartesianas, cilíndricas, esféricas, o en cualquier otro sistema
conveniente para el problema. Un problema con simetría esférica (como una
fuerza central) resulta mucho más simple en coordenadas esféricas que en
cartesianas. El objetivo de este capítulo es formalizar qué significa
``cambiar de coordenadas'' y cómo transforman, bajo dicho cambio, los
distintos objetos geométricos asociados a un vector: sus componentes, la
base coordenada, la base dual y la métrica.

\section{Coordinatización del espacio}

\begin{definicion}[Coordinatización]
Sea $M$ el espacio físico (por ejemplo $\mathbb{R}^3$). Una
\textbf{coordinatización} de $M$ es un mapeo biyectivo
\[
x : M \longrightarrow \mathbb{R}^3, \qquad p \longmapsto x(p) =
(x^1,x^2,x^3),
\]
que asigna a cada punto $p \in M$ una terna ordenada de números reales,
llamados las \textbf{coordenadas} de $p$.
\end{definicion}

\begin{ejemplo}
Las coordenadas cartesianas, cilíndricas y esféricas son tres
coordinatizaciones distintas de un mismo punto $p \in \mathbb{R}^3$:
\[
x_{\text{cart}}(p) = (x,y,z), \qquad
x_{\text{cil}}(p) = (\rho,\varphi,z), \qquad
x_{\text{esf}}(p) = (r,\varphi,\theta).
\]
\end{ejemplo}

\section{Transformación de coordenadas}

\begin{definicion}[Transformación de coordenadas]
Sean $x^i$ y $x'^i$ ($i=1,2,3$) dos coordinatizaciones distintas de un
mismo punto $p$. Una \textbf{transformación de coordenadas} es un par de
mapeos biyectivos e inversos entre sí,
\[
x'^i = x'^i(x^j), \qquad x^i = x^i(x'^j),
\]
que relacionan ambas coordinatizaciones.
\end{definicion}

\begin{ejemplo}[Cartesianas a esféricas]
La transformación de coordenadas cartesianas $(x,y,z)$ a esféricas
$(r,\varphi,\theta)$ está dada por
\[
x = r\sin\theta\cos\varphi, \qquad y = r\sin\theta\sin\varphi, \qquad
z = r\cos\theta,
\]
con transformación inversa
\[
r = \sqrt{x^2+y^2+z^2}, \qquad
\varphi = \arctan\!\left(\frac{y}{x}\right), \qquad
\theta = \arccos\!\left(\frac{z}{\sqrt{x^2+y^2+z^2}}\right).
\]
\end{ejemplo}

\subsection{El jacobiano de la transformación}

\begin{definicion}[Matriz jacobiana]
Dada una transformación de coordenadas $x'^i=x'^i(x^j)$, se define la
\textbf{matriz jacobiana} de la transformación como la matriz de
$3\times 3$
\[
J = \left[ \frac{\partial x'^i}{\partial x^j} \right]_{i,j=1,2,3} =
\begin{pmatrix}
\dfrac{\partial x'^1}{\partial x^1} & \dfrac{\partial x'^1}{\partial x^2} &
\dfrac{\partial x'^1}{\partial x^3} \\[8pt]
\dfrac{\partial x'^2}{\partial x^1} & \dfrac{\partial x'^2}{\partial x^2} &
\dfrac{\partial x'^2}{\partial x^3} \\[8pt]
\dfrac{\partial x'^3}{\partial x^1} & \dfrac{\partial x'^3}{\partial x^2} &
\dfrac{\partial x'^3}{\partial x^3}
\end{pmatrix},
\]
y el \textbf{jacobiano} de la transformación es su determinante,
$J = \det\!\left[\partial x'^i/\partial x^j\right]$.
\end{definicion}

\begin{observacion}
Para que la transformación sea invertible en un entorno de un punto se
requiere $J\neq 0$ en dicho punto (teorema de la función inversa). Los puntos donde $J=0$ son puntos singulares de la coordinatización; por
ejemplo, el origen $r=0$ en coordenadas esféricas, o el eje $z$
($\rho=0$) en coordenadas cilíndricas.
\end{observacion}

\begin{ejemplo}
Para la transformación de cartesianas a esféricas, el jacobiano es
\[
J = \det\!\left[\frac{\partial x'^i}{\partial x^j}\right] = r^2\sin\theta .
\]
El jacobiano se anula en $r=0$ y en $\theta=0,\pi$ (el eje $z$), que son
precisamente los puntos donde la coordinatización esférica es singular:
en el origen $\varphi$ y $\theta$ no están definidos, y sobre el eje $z$,
$\varphi$ no está definido.
\end{ejemplo}

\section{Vectores base}

\begin{definicion}[Base coordenada]
Dada una coordinatización $x^i$ de $\mathbb{R}^3$, la \textbf{base
coordenada} asociada es el conjunto de vectores
\[
\vec{e}_i = \frac{\partial \vec{r}}{\partial x^i}, \qquad i=1,2,3,
\]
donde $\vec{r}=\vec r(x^1,x^2,x^3)$ es el vector posición.
\end{definicion}

\begin{proposicionc}
Bajo una transformación de coordenadas $x^i = x^i(x'^j)$, la base
coordenada $\{\vec e_i\}_{i=1}^3$ transforma como
\[
\vec{e}\,'_i = \sum_{j=1}^{3} \frac{\partial x^j}{\partial x'^i}\, \vec e_j
.
\]
\end{proposicionc}

\begin{proof}
Por regla de la cadena,
\[
\vec e\,'_i = \frac{\partial \vec r}{\partial x'^i} =
\sum_{j=1}^{3} \frac{\partial x^j}{\partial x'^i}
\frac{\partial \vec r}{\partial x^j} =
\sum_{j=1}^{3} \frac{\partial x^j}{\partial x'^i}\, \vec e_j . \qedhere
\]
\end{proof}

De manera análoga, contrayendo la expresión anterior con
$\sum_i \partial x'^k/\partial x^i$ y usando
$\sum_i (\partial x'^k/\partial x^i)(\partial x^i/\partial x'^j) =
\delta^k_j$, se obtiene la relación inversa:

\begin{proposicionc}
Bajo la transformación inversa $x'^i=x'^i(x^j)$,
\[
\vec e_i = \sum_{j=1}^{3} \frac{\partial x'^j}{\partial x^i}\, \vec e\,'_j
.
\]
\end{proposicionc}

\begin{ejemplo}[Base esférica a partir de la cartesiana]
Sea $\{\vec e_1,\vec e_2,\vec e_3\} = \{\ihat,\jhat,\khat\}$ la base
cartesiana, y sea $\{\vec e\,'_1,\vec e\,'_2,\vec e\,'_3\} =
\{\vec e_r,\vec e_\varphi,\vec e_\theta\}$ la base coordenada esférica.
Usando $x^j=x^j(x'^i)=x^j(r,\varphi,\theta)$,
\[
\vec e\,'_1 = \vec e_r = \sum_{j=1}^3 \frac{\partial x^j}{\partial r}
\vec e_j = \frac{\partial x}{\partial r}\ihat + \frac{\partial y}{\partial
r}\jhat + \frac{\partial z}{\partial r}\khat ,
\]
lo que da
\[
\vec e_r = \sin\theta\cos\varphi\,\ihat + \sin\theta\sin\varphi\,\jhat +
\cos\theta\,\khat .
\]
Análogamente, derivando respecto de $\varphi$ y de $\theta$,
\[
\vec e_\varphi = r\sin\theta\,(-\sin\varphi\,\ihat + \cos\varphi\,\jhat),
\qquad
\vec e_\theta = r\big(\cos\varphi\cos\theta\,\ihat +
\sin\varphi\cos\theta\,\jhat - \sin\theta\,\khat\big).
\]

\medskip
Notar que $\vec e_r$ es unitario, pero $\vec e_\varphi$ y $\vec e_\theta$
no lo son: $\|\vec e_\varphi\| = r\sin\theta$ y $\|\vec e_\theta\|=r$. La
base coordenada $\{\vec e_r,\vec e_\varphi,\vec e_\theta\}$ es ortogonal
pero \textbf{no ortonormal}; ésta es la base ``natural'' asociada a las
coordenadas esféricas, distinta de la base ortonormal
$\{\hat r,\hat\varphi,\hat\theta\}$ usada habitualmente en física
elemental.

\medskip
Realizando el cálculo inverso, $x^i=x^i(x'^j)$, se obtienen los vectores
cartesianos en términos de la base esférica:
\[
\ihat = \frac{x}{\sqrt{x^2+y^2+z^2}}\,\vec e_r - \frac{y}{x^2+y^2}\,
\vec e_\varphi + \frac{xz}{\sqrt{x^2+y^2}\,(x^2+y^2+z^2)}\, \vec e_\theta ,
\]
y expresiones análogas para $\jhat$ y $\khat$; reescritas en coordenadas
esféricas por estética,
\[
\ihat = \cos\varphi\sin\theta\,\vec e_r - \frac{\sin\varphi}{r\sin\theta}\,
\vec e_\varphi + \frac{\cos\varphi\cos\theta}{r}\,\vec e_\theta, \quad
\jhat = \sin\varphi\sin\theta\,\vec e_r + \frac{\cos\varphi}{r\sin\theta}\,
\vec e_\varphi + \frac{\sin\varphi\cos\theta}{r}\,\vec e_\theta,
\]
\[
\khat = \cos\theta\,\vec e_r - \frac{\sin\theta}{r}\,\vec e_\theta .
\]
\end{ejemplo}

\section{Componentes de un vector: tensores contravariantes}

\begin{proposicionc}
\label{prop:transf-componentes}
Sea $\vec A$ un vector con componentes $A^i$ en la base $\{\vec e_i\}$.
Bajo la transformación de coordenadas $x^i=x^i(x'^j)$, las componentes de
$\vec A$ transforman como
\[
A'^i = \sum_{j=1}^{3} \frac{\partial x'^i}{\partial x^j}\, A^j .
\]
\end{proposicionc}

\begin{proof}
Como $\vec A$ no depende del sistema coordenado, $\vec A\big|_{O'} =
\vec A\big|_O$, es decir
\[
\sum_{i=1}^3 A'^i \vec e\,'_i = \sum_{j=1}^3 A^j \vec e_j =
\sum_{j=1}^3 A^j \sum_{i=1}^3 \frac{\partial x'^i}{\partial x^j}\vec e\,'_i
= \sum_{i=1}^3 \left(\sum_{j=1}^3 \frac{\partial x'^i}{\partial x^j}
A^j\right) \vec e\,'_i .
\]
Como la base $\{\vec e\,'_i\}$ es linealmente independiente, se concluye
\[
A'^i = \sum_{j=1}^3 \frac{\partial x'^i}{\partial x^j} A^j . \qedhere
\]
\end{proof}

Contrayendo con $\sum_i \partial x^k/\partial x'^i$ y usando la regla de
la cadena se obtiene la relación inversa:

\begin{proposicionc}
\[
A^k = \sum_{i=1}^{3} \frac{\partial x^k}{\partial x'^i}\, A'^i .
\]
\end{proposicionc}

\begin{ejemplo}
Sea $\vec A = A_x\ihat + A_y\jhat + A_z\khat = A_r\vec e_r +
A_\varphi\vec e_\varphi + A_\theta\vec e_\theta$. Usando la
Proposición~\ref{prop:transf-componentes} con la transformación a
esféricas se obtiene, tras simplificar,
\[
A_r = \cos\varphi\sin\theta\,A_x + \sin\varphi\sin\theta\,A_y +
\cos\theta\,A_z = \vec e_r\cdot\vec A ,
\]
\[
A_\varphi = -\frac{\sin\varphi}{r\sin\theta}A_x +
\frac{\cos\varphi}{r\sin\theta}A_y = \frac{1}{r^2\sin^2\theta}\,
\vec e_\varphi\cdot\vec A , \qquad
A_\theta = \frac{\cos\varphi\cos\theta}{r}A_x +
\frac{\sin\varphi\cos\theta}{r}A_y - \frac{\sin\theta}{r}A_z =
\frac{1}{r^2}\,\vec e_\theta\cdot\vec A .
\]
Notar que $A_\varphi$ y $A_\theta$ \textbf{no} son simplemente
$\vec e_\varphi\cdot\vec A$ ni $\vec e_\theta\cdot\vec A$, sino que llevan
un factor de normalización: esto ocurre porque $\{\vec e_r,\vec
e_\varphi,\vec e_\theta\}$ no es una base ortonormal (véase la sección
sobre la base dual, más adelante).

La transformación inversa da
\[
A_x = \frac{x}{\sqrt{x^2+y^2+z^2}}A_r - y A_\varphi +
\frac{xz}{\sqrt{x^2+y^2}}A_\theta, \qquad
A_y = \frac{y}{\sqrt{x^2+y^2+z^2}}A_r + x A_\varphi +
\frac{yz}{\sqrt{x^2+y^2}}A_\theta,
\]
\[
A_z = \frac{z}{\sqrt{x^2+y^2+z^2}}A_r - \sqrt{x^2+y^2}\,A_\theta .
\]
\end{ejemplo}

\begin{observacion}
De la discusión anterior se concluye que:
\begin{itemize}
\item Las cantidades con el índice \textbf{arriba} transforman con la
matriz $\partial x'^i/\partial x^j$.
\item Las cantidades con el índice \textbf{abajo} transforman con la
matriz $\partial x^i/\partial x'^j$.
\end{itemize}
En la construcción de un vector $\vec A = \sum_i A^i \vec e_i$ se
necesitan las cantidades $\{\vec e_i\}$ (índice abajo) y $\{A^i\}$
(índice arriba): se contrae un índice arriba con uno abajo. Esta es la
razón estructural por la que los vectores no dependen del sistema
coordenado: las dos leyes de transformación (una y su inversa) se
cancelan exactamente.
\end{observacion}

\begin{definicion}[Índices contraídos]
Un índice arriba y un índice abajo se dicen \textbf{contraídos} si están
bajo una suma sobre dicho índice.
\end{definicion}

\begin{definicion}[Tensor contravariante de orden uno]
La cantidad $T^i$ se dice un \textbf{tensor contravariante de orden uno},
o \textbf{tensor de orden $(1,0)$}, si bajo una transformación de
coordenadas $x^i=x^i(x'^j)$ transforma como
\[
T'^i = \sum_{j=1}^{3} \frac{\partial x'^i}{\partial x^j}\, T^j .
\]
\end{definicion}

Las componentes $A^i$ de un vector arbitrario $\vec A$ son, por lo tanto,
un tensor de orden $(1,0)$.

\section{Vectores duales y tensores covariantes}

\begin{definicion}[Base dual]
La \textbf{base dual} $\{\vec E^i\}_{i=1}^3$ asociada a la base
coordenada $\{\vec e_i\}_{i=1}^3$ es el conjunto de (co)vectores definido
por la relación invariante
\[
E^i(\vec e_j) = \delta^i_j ,
\]
que se cumple en cualquier sistema de coordenadas:
$E^i(e_j)\big|_O = E^i(e_j)\big|_{O'} = \delta^i_j$.
\end{definicion}

\begin{observacion}
Los vectores duales pertenecen al espacio dual $V^*(n,K)$; un vector dual
arbitrario se escribe $\vec f = \sum_i f_i \vec E^i$, con componentes con
el índice abajo. Como $\vec f$ tampoco depende del sistema coordenado
($\vec f\big|_{O'}=\vec f\big|_O$), la base dual debe poseer su propia
ley de transformación, distinta de la de la base coordenada, ya que
$\vec E^i$ lleva el índice arriba.
\end{observacion}

\begin{proposicionc}
\label{prop:transf-base-dual}
Bajo la transformación de coordenadas $x^i=x^i(x'^j)$, la base dual
transforma como
\[
E'^i = \sum_{j=1}^{3} \frac{\partial x'^i}{\partial x^j}\, E^j .
\]
\end{proposicionc}

\begin{proof}
Partimos de la relación invariante $E^i(e_j)\big|_{O'} = \delta^i_j$ y
usamos la ley de transformación de $\vec e_j$ (transformación inversa
$x^j=x^j(x'^k)$):
\[
E'^i(e'_j) = \left(\sum_{k=1}^3 \frac{\partial x'^i}{\partial x^k} E^k
\right)\!\left(\sum_{l=1}^3 \frac{\partial x^l}{\partial x'^j} e_l\right)
= \sum_{k=1}^3\sum_{l=1}^3 \frac{\partial x'^i}{\partial x^k}
\frac{\partial x^l}{\partial x'^j}\, E^k(e_l)
= \sum_{k=1}^3 \frac{\partial x'^i}{\partial x^k}\frac{\partial
x^k}{\partial x'^j} = \delta^i_j , \qedhere
\]
lo cual confirma que la ley propuesta es consistente con la definición
invariante de la base dual.
\end{proof}

Contrayendo con $\sum_i \partial x^k/\partial x'^i$ se obtiene la
transformación inversa:
\[
E^k = \sum_{i=1}^{3} \frac{\partial x^k}{\partial x'^i}\, E'^i .
\]

\begin{proposicionc}
\label{prop:transf-componentes-dual}
La componente $f_i$ de un vector dual $\vec f$ transforma, bajo
$x'^i=x'^i(x^j)$, como
\[
f'_i = \sum_{j=1}^{3} \frac{\partial x^j}{\partial x'^i}\, f_j .
\]
\end{proposicionc}

\begin{proof}
Como $\vec f\big|_O = \vec f\big|_{O'}$,
\[
\sum_{j=1}^3 f_j E^j = \sum_{i=1}^3 f'_i E'^i = \sum_{i=1}^3 f'_i
\sum_{j=1}^3 \frac{\partial x^j}{\partial x'^i} E^j ,
\]
y por independencia lineal de $\{\vec E^j\}$,
\[
f'_i = \sum_{j=1}^3 \frac{\partial x^j}{\partial x'^i} f_j . \qedhere
\]
\end{proof}

\begin{definicion}[Tensor covariante de orden uno]
La cantidad $T_i$ se dice un \textbf{tensor covariante de orden uno}, o
\textbf{tensor de orden $(0,1)$}, si bajo la transformación de
coordenadas $x^i=x^i(x'^j)$ transforma como
\[
T'_i = \sum_{j=1}^{3} \frac{\partial x^j}{\partial x'^i}\, T_j .
\]
\end{definicion}

\begin{observacion}
El término \textbf{covarianza} (invariancia de forma) de un vector
significa que el vector no cambia de forma bajo una transformación de
coordenadas: no cambia su dirección, sentido ni magnitud. Sin embargo,
bajo un cambio de coordenadas cambian las bases coordenadas, por lo que
el valor numérico de sus componentes sí cambia.
\end{observacion}

\section{La métrica}

\begin{definicion}[Métrica]
La \textbf{métrica} asociada a una base coordenada $\{\vec e_i\}_{i=1}^3$
es
\[
g_{ij} = \vec e_i \cdot \vec e_j .
\]
\end{definicion}

\begin{ejemplo}
En coordenadas cartesianas, $g_{ij}=\delta_{ij}$ (la \textbf{métrica de
Euclides}). En coordenadas esféricas,
\[
g_{ij} = \begin{pmatrix} 1 & 0 & 0 \\ 0 & r^2\sin^2\theta & 0 \\ 0 & 0 &
r^2 \end{pmatrix},
\]
consistente con $\|\vec e_r\|^2=1$, $\|\vec e_\varphi\|^2 = r^2\sin^2
\theta$, $\|\vec e_\theta\|^2=r^2$ obtenidas anteriormente.
\end{ejemplo}

\begin{proposicionc}
\label{prop:transf-metrica}
La métrica $g_{ij}(x^k)$ bajo una transformación de coordenadas
$x^i=x^i(x'^j)$ transforma como
\[
g'_{ij} = \sum_{k=1}^{3}\sum_{l=1}^{3} \frac{\partial x^k}{\partial x'^i}
\frac{\partial x^l}{\partial x'^j}\, g_{kl} .
\]
\end{proposicionc}

\begin{proof}
Usando la ley de transformación de la base coordenada,
\[
g'_{ij} = \vec e\,'_i \cdot \vec e\,'_j = \left(\sum_{k=1}^3
\frac{\partial x^k}{\partial x'^i}\vec e_k\right)\!\cdot\!\left(
\sum_{l=1}^3 \frac{\partial x^l}{\partial x'^j}\vec e_l\right)
= \sum_{k=1}^3\sum_{l=1}^3 \frac{\partial x^k}{\partial x'^i}
\frac{\partial x^l}{\partial x'^j}\, \vec e_k\cdot\vec e_l =
\sum_{k=1}^3\sum_{l=1}^3 \frac{\partial x^k}{\partial x'^i}\frac{\partial
x^l}{\partial x'^j} g_{kl} . \qedhere
\]
\end{proof}

\begin{definicion}[Tensor covariante de orden dos]
La cantidad $T_{ij}$ se dice un tensor covariante de orden $2$, o tensor
de orden $(0,2)$, si transforma como
\[
T'_{ij} = \sum_{k=1}^{3}\sum_{l=1}^{3} \frac{\partial x^k}{\partial
x'^i}\frac{\partial x^l}{\partial x'^j}\, T_{kl} .
\]
\end{definicion}

Por lo tanto, la métrica $g_{ij}$ es un tensor de orden $(0,2)$.

\begin{ejemplo}
Recuperemos la métrica esférica a partir de la métrica cartesiana
$g_{ij}=\delta_{ij}$ usando la Proposición~\ref{prop:transf-metrica}:
\[
g'_{ij} = \sum_{k=1}^3 \frac{\partial x^k}{\partial x'^i}\frac{\partial
x^k}{\partial x'^j} = \frac{\partial x}{\partial x'^i}\frac{\partial
x}{\partial x'^j} + \frac{\partial y}{\partial x'^i}\frac{\partial
y}{\partial x'^j} + \frac{\partial z}{\partial x'^i}\frac{\partial
z}{\partial x'^j} .
\]
Sustituyendo $x=r\sin\theta\cos\varphi$, etc., se recupera
\[
g'_{ij} = \begin{pmatrix} 1 & 0 & 0 \\ 0 & r^2\sin^2\theta & 0 \\ 0 & 0 &
r^2 \end{pmatrix}.
\]
\end{ejemplo}

\subsection{El elemento de línea}

\begin{proposicionc}
El elemento de línea $ds$, definido por $ds^2 = \sum_{i,j} g_{ij}\,
dx^i\,dx^j$, es un invariante bajo transformaciones de coordenadas:
$ds' = ds$.
\end{proposicionc}

\begin{proof}
En el sistema $O'x'^i$,
\[
ds'^2 = \sum_{i=1}^3\sum_{j=1}^3 g'_{ij}\, dx'^i\, dx'^j .
\]
Bajo $x'^i=x'^i(x^j)$, el diferencial $dx'^i$ transforma como
$dx'^i = \sum_j (\partial x'^i/\partial x^j)\, dx^j$, es decir, $dx^i$ es
un tensor de orden $(1,0)$. Sustituyendo esto y la ley de transformación
de $g'_{ij}$ (Proposición~\ref{prop:transf-metrica}),
\[
ds'^2 = \sum_{i,j,k,l} \frac{\partial x^k}{\partial x'^i}\frac{\partial
x^l}{\partial x'^j}\, g_{kl} \sum_{n} \frac{\partial x'^i}{\partial x^n}
dx^n \sum_{m}\frac{\partial x'^j}{\partial x^m} dx^m
= \sum_{k,l,n,m} \left(\sum_i \frac{\partial x^k}{\partial x'^i}
\frac{\partial x'^i}{\partial x^n}\right)\!\!\left(\sum_j \frac{\partial
x^l}{\partial x'^j}\frac{\partial x'^j}{\partial x^m}\right) g_{kl}\,
dx^n dx^m .
\]
Ambos paréntesis son deltas de Kronecker, $\delta^k_n$ y $\delta^l_m$, de
modo que
\[
ds'^2 = \sum_{k,l} g_{kl}\, dx^k dx^l = ds^2 \quad\Longrightarrow\quad
ds' = ds . \qedhere
\]
\end{proof}

\begin{observacion}
El elemento de línea no es un vector sino un \emph{escalar}: un número
(con dimensión) cuyo valor no cambia bajo un cambio de coordenadas. En
particular, el largo de una trayectoria tiene el mismo valor en todos los
sistemas coordenados.
\end{observacion}

\section{Invariantes: producto escalar, magnitud y ángulo}

\begin{proposicionc}
El producto escalar $\vec A\cdot\vec B = \sum_{i,j} g_{ij} A^i B^j$ entre
dos vectores cualesquiera es invariante bajo transformación de
coordenadas.
\end{proposicionc}

\begin{proof}
Análogo a la demostración de la invariancia del elemento de línea,
reemplazando $dx^i$ por $A^i$ y $dx^j$ por $B^j$ (ambos tensores de orden
$(1,0)$):
\[
\vec A\cdot\vec B\big|_{O'} = \sum_{i,j} g'_{ij} A'^i B'^j = \sum_{n,m}
g_{nm} A^n B^m = \vec A\cdot\vec B\big|_O . \qedhere
\]
\end{proof}

El producto escalar es un número real, no un vector; por lo tanto su
valor numérico no cambia bajo una transformación de coordenadas. De este
resultado se derivan dos corolarios inmediatos.

\begin{corolario}
La magnitud $\|\vec A\| = \sqrt{\vec A\cdot\vec A}$ de un vector es
invariante bajo transformaciones de coordenadas.
\end{corolario}

\begin{proof}
Tomando $\vec B=\vec A$ en la proposición anterior,
$\vec A\cdot\vec A\big|_{O'} = \vec A\cdot\vec A\big|_O$, es decir
$\|\vec A\|^2\big|_{O'} = \|\vec A\|^2\big|_O$, luego
$\|\vec A\|\big|_{O'} = \|\vec A\|\big|_O$. \qedhere
\end{proof}

\begin{corolario}
El ángulo entre dos vectores es invariante bajo transformación de
coordenadas.
\end{corolario}

\begin{proof}
El ángulo $\theta$ entre $\vec A$ y $\vec B$ está dado por
$\cos\theta = (\vec A\cdot\vec B)/(\|\vec A\|\|\vec B\|)$. Como el
numerador y ambos factores del denominador son invariantes,
$\cos\theta\big|_{O'} = \cos\theta\big|_O$, y por lo tanto
$\theta\big|_{O'} = \theta\big|_O$. \qedhere
\end{proof}

\begin{observacion}[Definición formal vía el espacio dual]
Formalmente, el producto escalar entre un vector dual $\vec f = \sum_i
f_i \vec E^i \in V^*(n,K)$ y un vector $\vec v = \sum_i v^i \vec e_i \in
V(n,K)$ se define como el emparejamiento natural
\[
f(v) = \sum_{i=1}^{3} f_i\, v^i .
\]
\end{observacion}

\begin{proposicionc}
El emparejamiento $f(v) = \sum_i f_i v^i$ es invariante bajo la
transformación de coordenadas $x'^i=x'^i(x^j)$.
\end{proposicionc}

\begin{proof}
Las cantidades $v^i$ y $f_i$ son tensores de orden $(1,0)$ y $(0,1)$
respectivamente, es decir $v'^i = \sum_j (\partial x'^i/\partial x^j)
v^j$ y $f'_i = \sum_j (\partial x^j/\partial x'^i) f_j$. Entonces
\[
f'(v') = \sum_{i=1}^3 f'_i v'^i = \sum_{i,j,k}\frac{\partial
x^j}{\partial x'^i} f_j\, \frac{\partial x'^i}{\partial x^k} v^k =
\sum_{j,k}\left(\sum_i \frac{\partial x^j}{\partial x'^i}\frac{\partial
x'^i}{\partial x^k}\right) f_j v^k = \sum_{j,k} \delta^j_k\, f_j v^k =
\sum_{j} f_j v^j = f(v). \qedhere
\]
\end{proof}

\newpage
\section{La métrica inversa}

\begin{proposicionc}
La métrica inversa $g^{ij}$ está dada en términos de la base dual
$\{\vec E^i\}_{i=1}^3$ por
\[
g^{ij} = \vec E^i \cdot \vec E^j .
\]
\end{proposicionc}

\begin{proof}
La base dual se relaciona con la base coordenada a través de la métrica
inversa, $\vec E^i = \sum_j g^{ij} \vec e_j$, de modo que se cumpla la
definición $E^i(e_j) = \vec E^i\cdot\vec e_j = \sum_k g^{ik}\vec
e_k\cdot\vec e_j = \sum_k g^{ik} g_{kj} = \delta^i_j$. Entonces
\[
\vec E^i\cdot\vec E^j = \sum_{k} g^{ik}\vec e_k \cdot \sum_{l} g^{jl}\vec
e_l = \sum_{k,l} g^{ik} g^{jl}\, g_{kl} = \sum_k g^{ik}\left(\sum_l
g^{jl}g_{kl}\right) = \sum_k g^{ik}\, \delta^j_k = g^{ij} . \qedhere
\]
\end{proof}

\begin{ejemplo}
En coordenadas esféricas, usando $\vec E^1=\vec e_r$, $\vec E^2 = \vec
e_\varphi/(r^2\sin^2\theta)$, $\vec E^3 = \vec e_\theta/r^2$ (obtenidos de
$\vec E^i = \sum_j g^{ij}\vec e_j$ con $g^{ij}=\text{diag}(1,\,
1/(r^2\sin^2\theta),\,1/r^2)$), se verifica directamente
\[
g^{11} = \vec E^1\cdot\vec E^1 = \vec e_r\cdot\vec e_r = 1, \qquad
g^{22} = \vec E^2\cdot\vec E^2 = \frac{\vec e_\varphi\cdot\vec
e_\varphi}{r^4\sin^4\theta} = \frac{1}{r^2\sin^2\theta}, \qquad
g^{33} = \vec E^3\cdot\vec E^3 = \frac{\vec e_\theta\cdot\vec
e_\theta}{r^4} = \frac{1}{r^2},
\]
y los términos cruzados $g^{12}=g^{13}=g^{23}=0$ por ortogonalidad de la
base esférica. Es decir,
\[
g^{ij} = \begin{pmatrix} 1 & 0 & 0 \\ 0 & \dfrac{1}{r^2\sin^2\theta} & 0
\\ 0 & 0 & \dfrac{1}{r^2} \end{pmatrix},
\]
el resultado conocido.
\end{ejemplo}

\begin{proposicionc}
\label{prop:transf-metrica-inversa}
La métrica inversa $g^{ij}$ bajo una transformación de coordenadas
$x'^i=x'^i(x^j)$ transforma como
\[
g'^{ij} = \sum_{k=1}^{3}\sum_{l=1}^{3} \frac{\partial x'^i}{\partial
x^k}\frac{\partial x'^j}{\partial x^l}\, g^{kl} .
\]
\end{proposicionc}

\begin{proof}
Escribiendo $g'^{ij} = \vec E'^i\cdot\vec E'^j$ y usando la ley de
transformación de la base dual (Proposición~\ref{prop:transf-base-dual}),
\[
g'^{ij} = \left(\sum_k \frac{\partial x'^i}{\partial x^k}\vec
E^k\right)\!\cdot\!\left(\sum_l \frac{\partial x'^j}{\partial x^l}\vec
E^l\right) = \sum_{k,l} \frac{\partial x'^i}{\partial x^k}\frac{\partial
x'^j}{\partial x^l}\, \vec E^k\cdot\vec E^l = \sum_{k,l} \frac{\partial
x'^i}{\partial x^k}\frac{\partial x'^j}{\partial x^l}\, g^{kl} .
\qedhere
\]
\end{proof}

\begin{ejemplo}
Verifiquemos la componente $g'^{11}$ de la métrica esférica a partir de
$g^{kl}=\delta^{kl}$ (cartesiana) usando la
Proposición~\ref{prop:transf-metrica-inversa}:
\[
g'^{11} = \left(\frac{\partial r}{\partial x}\right)^2 +
\left(\frac{\partial r}{\partial y}\right)^2 + \left(\frac{\partial
r}{\partial z}\right)^2 = \frac{x^2}{r^2} + \frac{y^2}{r^2} +
\frac{z^2}{r^2} = 1 ,
\]
consistente con el resultado directo. Repitiendo el cálculo para las
demás componentes se recupera $g^{ij} =
\text{diag}(1,\,1/(r^2\sin^2\theta),\,1/r^2)$.
\end{ejemplo}

\begin{definicion}[Tensor contravariante de orden dos]
La cantidad $T^{ij}$ se dice un tensor contravariante de orden $2$, o
tensor de orden $(2,0)$, si bajo $x'^i=x'^i(x^j)$ transforma como
\[
T'^{ij} = \sum_{k=1}^{3}\sum_{l=1}^{3} \frac{\partial x'^i}{\partial
x^k}\frac{\partial x'^j}{\partial x^l}\, T^{kl} .
\]
\end{definicion}

Por lo tanto, la métrica inversa $g^{ij}$ es un tensor de orden $(2,0)$.

\section{Los tres símbolos delta}

Conviene distinguir con cuidado tres cantidades que, a primera vista,
parecen representar todas a la matriz identidad: $\delta_{ij}$,
$\delta^{ij}$ y $\delta^i_j$.

\begin{proposicionc}
De las tres cantidades $\delta_{ij}$, $\delta^{ij}$ y $\delta^i_j$, solo
$\delta^i_j$ es invariante bajo transformaciones de coordenadas.
\end{proposicionc}

\begin{proof}
Como $\delta^k_l$ es un tensor de orden $(1,1)$ (transforma con un
índice arriba y uno abajo, cada uno con su matriz correspondiente),
\[
\delta'^i_j = \sum_{k=1}^3\sum_{l=1}^3 \frac{\partial x'^i}{\partial
x^k}\frac{\partial x^l}{\partial x'^j}\, \delta^k_l = \sum_{k=1}^3
\frac{\partial x'^i}{\partial x^k}\frac{\partial x^k}{\partial x'^j} =
\frac{\partial x'^i}{\partial x'^j} = \delta^i_j . \qedhere
\]
\end{proof}

En cambio, $\delta_{ij}$ y $\delta^{ij}$ transforman como tensores de
orden $(0,2)$ y $(2,0)$ respectivamente, y por lo tanto \textbf{no} son
invariantes: cambian su forma funcional bajo un cambio de coordenadas
(por ejemplo, al pasar a esféricas dejan de ser la matriz identidad).

\begin{observacion}
En resumen:
\begin{itemize}
\item $\delta^i_j$ es la matriz identidad (invariante, tensor de orden
$(1,1)$).
\item $\delta_{ij}$ es la métrica asociada a las coordenadas cartesianas,
llamada \textbf{métrica de Euclides} (tensor de orden $(0,2)$).
\item $\delta^{ij}$ es la métrica inversa de la métrica de Euclides
(tensor de orden $(2,0)$).
\end{itemize}
\end{observacion}
